% !TeX spellcheck = en_US
\documentclass[aspectratio=1610]{beamer}
\usetheme{Madrid}
\setbeamercovered{transparent}
\setbeamertemplate{navigation symbols}{}
\setbeamertemplate{itemize items}[circle]
\setbeamertemplate{enumerate items}[default]

\makeatletter
\define@key{beamerframe}{wide}[10pt]{%
	\def\beamer@cramped{\itemsep #1\topsep0.5pt\relax}}
\makeatother
	

\usepackage[english]{babel}
\usepackage[utf8]{inputenc}
\usepackage[T1]{fontenc}
\usepackage{array}
\usepackage{xcolor}
\usepackage{colortbl}
\usepackage{booktabs}

\title[Fall - Lecture 5] % (optional, use only with long paper titles)
{Intermediate Macroeconomics: Lecture  5 (Ch. 4)}
\subtitle{ECON 311 -- Intermediate Macroeconomics}

\author{Muhebullah Karimzada}
\institute{School of Economics \\\ University of Northern British Columbia}
\date{Fall 2026}


% ---------- UNBC COLORS ----------
\usepackage{xcolor}
\definecolor{UNBCGreen}{RGB}{0,102,51}
\definecolor{UNBCDark}{RGB}{0,74,37}
\definecolor{UNBCGold}{RGB}{189,155,96}
\definecolor{SoftCream}{RGB}{249,247,242}
\definecolor{SoftGray}{RGB}{244,246,248}
\definecolor{AccentBlue}{RGB}{41,98,255}
\definecolor{AccentOrange}{RGB}{230,126,34}
\definecolor{AccentTeal}{RGB}{0,150,136}
\definecolor{AccentRed}{RGB}{192,57,43}
\setbeamercolor{normal text}{fg=black,bg=white}
\setbeamercolor{structure}{fg=UNBCDark}
\setbeamercolor{title}{fg=white,bg=UNBCDark}
\setbeamercolor{subtitle}{fg=UNBCGold}
\setbeamercolor{frametitle}{fg=white,bg=UNBCDark}
\setbeamercolor{block title}{fg=white,bg=UNBCDark}
\setbeamercolor{block body}{fg=black,bg=SoftGray}
\setbeamercolor{block title alerted}{fg=white,bg=AccentRed}
\setbeamercolor{block body alerted}{fg=black,bg=SoftGray}
\setbeamercolor{item}{fg=UNBCDark}
\setbeamercolor{section in toc}{fg=UNBCDark}
\setbeamercolor{alerted text}{fg=AccentRed}
\setbeamercolor{author in head/foot}{fg=white,bg=UNBCDark}
\setbeamercolor{date in head/foot}{fg=white,bg=UNBCDark}
% ---------- UNBC BRANDING ----------
\usepackage{graphicx}
\usepackage{lmodern}
\usepackage{tcolorbox}
\usepackage{tikz}
\usetikzlibrary{positioning,calc}
\setbeamertemplate{headline}{}
\setbeamertemplate{footline}{%
  \leavevmode%
  \hbox{%
  \begin{beamercolorbox}[wd=.78\paperwidth,ht=2.8ex,dp=1.2ex,leftskip=1em]{author in head/foot}%
  \color{white}\textbf{ECON 311} \hspace{0.5em} \textcolor{UNBCGold}{|} \hspace{0.5em} Intermediate Macroeconomics%
  \end{beamercolorbox}%
  \begin{beamercolorbox}[wd=.22\paperwidth,ht=2.8ex,dp=1.2ex,rightskip=1em plus1fil]{date in head/foot}%
  \color{white}\insertframenumber/\inserttotalframenumber%
  \end{beamercolorbox}}%
}

\newtcolorbox{mybox}[2][]{
  colback=#2!8,
  colframe=#2,
  boxrule=0.8pt,
  arc=2mm,
  left=2mm,right=2mm,top=1.5mm,bottom=1.5mm,
  #1
}

\newcommand{\topbanner}[1]{%
\begin{tikzpicture}[remember picture,overlay]
\fill[UNBCGold] ([yshift=-0.65cm]current page.north west) rectangle ([yshift=-0.48cm]current page.north east);
\node[anchor=north west,font=\small\bfseries,text=UNBCDark] at ([xshift=0.35cm,yshift=-0.78cm]current page.north west) {#1};
\end{tikzpicture}%
}
\begin{document}

\begin{frame}[plain]
  \titlepage
\end{frame}

\begin{frame}[wide]{Firm's Problem}
	\begin{itemize}
		\item 	Firms decide how many workers and machines to use in production
		\[Y = F(K,L)\]
		\begin{itemize}
			\item  Choose K and L to maximize Profits = Revenue - Cost
		\end{itemize}
	\end{itemize}
\end{frame}

\begin{frame}[wide]{Firm's Problem}
	\begin{itemize}
		\item Profit is defined as total revenue minus total costs:
		\[\max_{K,L} \pi = F(K,L) - rK - wL\]
		where $\pi$ is the profit, $r$ is the rental rate of capital, and $w$ is the wage rate
		
		\begin{itemize}
			\item The term 	$rK$ is the price per unit of capital multiplied by the number of capital inputs used (total cost of all capital)
			\item wL is the price of labour multiplied by the number of labour inputs used (total cost of all labour)
		\end{itemize}
	\end{itemize}
\end{frame}

\begin{frame}[wide]{Firm's Problem}
	\begin{itemize}
		\item Assuming a perfectly competitive market
		\item Firms take rental rate $r$ and wage rate $w$ as given
		\vspace{2mm}
		\begin{itemize}
			\item No firm has control to change input prices
			\item No matter how much or how little L or K a firm uses, the prices won’t change
		\end{itemize}
		\item Input prices will change only as the result of an exogenous force in the market for K or L
		\item  For simplicity, the price of the output is normalized to one
	\end{itemize}
\end{frame}

\begin{frame}[wide]{Firm's Problem}
	\begin{itemize}
		\item The rental rate and wage rate are taken as given under perfect competition
		\begin{itemize}
			\item Hire capital until $MPK = r$
			\item Hire labour until $MPL = w$
			\item If $MPL > w$, the amount of output produced by an additional worker is greater than w
			\item Profits increase, this will continue until the cost of hiring a worker is equal to the amount of extra output the worker produces
			\item The same reasoning holds for capital
		\end{itemize}
		
	\end{itemize}
\end{frame}

\begin{frame}[wide]{Marginal Products}
	\begin{itemize}
		\item Using the Codd-Douglas production function: $Y=AK^{1/3} L^{2/3}$
		\item 	The marginal product of labour (MPL) is the additional output that is produced when one unit of labour is added, holding all other inputs constant:
		\[MPL = \frac{2}{3}A \left(\frac{K}{L}\right)^{1/3} = \left(\frac{2}{3}\right) \left(\frac{Y}{L}\right)\]
		\item The marginal product of capital (MPK) is the additional output that is produced when one unit of capital is added, holding all other inputs constant:
		\[MPK = \frac{1}{3} A \left(\frac{L}{K}\right)^{2/3} = \left(\frac{1}{3}\right)\left(\frac{Y}{K}\right)\]
		\item In a Cobb-Douglas production function, the marginal product of an input is equal to the product of the factor’s exponent times the average amount that each unit of the factor produces
		
	\end{itemize}
\end{frame}

\begin{frame}[wide]{Diminishing Marginal Product of Capital in Production}
	\begin{itemize}
		\item When labour input is held constant, Each additional unit of K increases output by less and less
	\end{itemize}
	\begin{figure}
		\centering
		\includegraphics[width=0.6\linewidth]{"Figures/Diminishing Marginal Product of Capital in Production"}
	\end{figure}

\end{frame}

\begin{frame}[wide]{Diminishing Returns}
	\begin{itemize}
		\item Formally, 
		\[\frac{\partial^2 F  }{\partial K^2}<0 \text{ and } \frac{\partial^2 F }{\partial L^2} <0\]		
		\item For example, suppose we have five computers
		\begin{itemize}
			\item  The first five people we hire are very productive
			\item But as we add more and more workers, there is less for them to do with only five computers among them
		\end{itemize}
	\end{itemize}
\end{frame}

\begin{frame}[wide]{Solving the Model: General Equilibrium}
	\begin{itemize}
		\item In general equilibrium, we solve for not only one signle market but both market of labour and capital
		\item Assume Number of perfectly competitive firms = 1, we could have more firms if we wished and nothing would change since our model features constant returns to scale
		\item Demand for capital (labour) by our representative firm = aggregate quantity of capital  (labour) in the economy		
	\end{itemize}
\end{frame}

\begin{frame}[wide]{Solving the Model: General Equilibrium}
	\begin{itemize}
		\item Five endogenous variables
		\begin{itemize}
			\item Output (Y)
			\item The amount of capital (K)
			\item The amount of labour (L)
			\item The wage (w)
			\item The rental price of capital (r)
		\end{itemize}
	\end{itemize}
\end{frame}

\begin{frame}[wide]{Solving the Model: General Equilibrium}
	\begin{itemize}
		\item Five equations
		\begin{itemize}
			\item The production function
			\[Y=AK^{1/3}L^{2/3}\]
			\item The rule for hiring capital
			\[\frac{1}{3}\frac{Y}{K} = r\]
			\item The rule for hiring labour
			\[\frac{2}{3}\frac{Y}{L} = w\]
			\item Supply equals the demand for labour
			\[L = \bar{L}\]
			\item Supply equals the demand for capital
			\[K = \bar{K}\]
		\end{itemize}
	\item The supply of labour and capital are assumed to be perfectly inelastic, the economy will supply $\bar{L}$ around of labour and $\bar{K}$ amount of capital regardless of the price
	\end{itemize}
\end{frame}

\begin{frame}[wide]{Solving the Model: Solution}
	\begin{itemize}
		\item 	A solution to the model
		\begin{itemize}
			\item A new set of equations that express the five unknowns in terms of the parameters and exogenous variables
			\item There should be five of them as we have five endogenous variables
			\item Endogenous variables are on the left side, and the parameters and exogenous variables that describe market behavior in equilibrium are on the right side.
			
		\end{itemize}
		\item General equilibrium
		\begin{itemize}
			\item Solution to the model when more than a single market clears
		\end{itemize}
		
	\end{itemize}
\end{frame}

\begin{frame}[wide]{Supply and Demand in the Capital and labour Markets}
	\begin{itemize}
		\item 	The solution easiest to see graphically (supply and demand for K and L)
		\item The supplies of K and L exogenously given (vertical lines)
		\item The demand curves for capital and labour are based on the hiring rules
		\begin{itemize}
			\item They trace out exactly the marginal product schedules for K and L
		\end{itemize}
	\end{itemize}
	\begin{figure}
		\centering
		\includegraphics[width=0.65\linewidth]{"Figures/Supply and Demand in the Capital and labor Markets"}
	\end{figure}
	
\end{frame}

\begin{frame}[wide]{In This Model…}
	\begin{itemize}
		\item 	The solution implies that
		\begin{itemize}
			\item Firms employ all the supplied capital and labour in the economy
			\item The production function is evaluated with the given supply of inputs: 
			\[Y^* = F(\bar{K},\bar{L}) = \bar{A}\bar{K}^\alpha \bar{L}^{1-\alpha}\]
			\item $w = MPL$ evaluated at the equilibrium values of Y, K, and L
			\item $r = MPK$ evaluated at the equilibrium values of Y, K, and L
		\end{itemize}
		
	\end{itemize}
\end{frame}

\begin{frame}[wide]{Interpreting the Solution}
	\begin{itemize}
		\item The equilibrium wage is proportional to output per worker: 	Output per worker = (Y/L)
		\item The equilibrium rental rate is proportional to output per capital:	Output per capital = (Y/K)
		\item In the United States, empirical evidence shows:
		\[\frac{w^*L^*}{Y^*} =\frac{2}{3}  \text{  and  } \frac{r^*K^*}{Y^*} =\frac{1}{3}\]
		\item The factor shares of the payments are equal to the exponents on the inputs in the Cobb-Douglas function
	\end{itemize}
\end{frame}

\begin{frame}[wide]{Equilibrium}
	All income is paid to capital or labour
	\begin{itemize}
		\item Verifies the assumption of perfect competition
		\item One of the implications of perfect competition is zero profits
			\[\pi = Y^* - r^*K^* - w^*L^* = 0\]
		\item Also verifies that production equals spending equals income
			\[w^*L^* + r^* K^* = Y^*\]
			\[\text{Income} = \text{production}\]
		\item This is one of the fundamental relationships of national income accounting.
	\end{itemize}
\end{frame}

\begin{frame}[wide]{Analyzing the Production Model}
	Development accounting:
	\begin{itemize}
		\item The use of a model to explain differences in incomes across countries
			\[y^* = \bar{A} \bar{k}^{1/3}\]
		\item For now, setting the productivity parameter = 1
			\[y^* = \bar{k}^{1/3}\]
		
	\end{itemize}
\end{frame}

\begin{frame}[wide]{The Empirical Fit of the Production Function }
	\begin{itemize}
		\item By using different countries' GDP per capita and capital per capita, we can check the empirical fit of the production function
		\item If the productivity parameter is 1, the model overpredicts GDP per capita
		\begin{itemize}
			\item Assuming that the productivity parameter is constant across countries, and does not vary over time
		\end{itemize}
		\item Capital per person varies tremendously across countries ranging from a low in Burundi of less than 1 percent of the U.S. value to a high in Switzerland of about 30 percent greater

	\end{itemize}
\end{frame}

\begin{frame}[wide]{The Model’s Prediction for Per Capita GDP (United States = 1)}
	\begin{itemize}
		\item 	The predicted value is simply $y=(0.93)^{1/3} = 0.976$ for Japan
	\end{itemize}
	\begin{figure}
		\centering
		\includegraphics[width=0.6\linewidth]{"Figures/The Model’s Prediction for Per Capita GDP"}
	\end{figure}
	
\end{frame}

\begin{frame}[wide]{The Empirical Fit of the Production Function }
	\begin{itemize}
		\item 	With no difference in productivity ($\bar{A} = 1$), the model predicts smaller differences in income across countries than we observe
		\item Diminishing returns to capital implies that
		\begin{itemize}
			\item countries with low K will have a high MPK
			\item countries with a lot of K will have a low MPK, and cannot raise GDP per capita by much through more capital accumulation
			
		\end{itemize}
		\item Differences in capital per person translate into much smaller differences in GDP because of diminishing MPK
		\item Therefore, gaps in income between rich and poor countries are predicted to be smaller than they are
	\end{itemize}
\end{frame}

\begin{frame}[wide]{Predicted Per Capita GDP in the Production Model}
	\begin{figure}
		\centering
		\includegraphics[width=0.6\linewidth]{"Figures/Predicted Per Capita GDP in the Production Model"}
	\end{figure}
	
\end{frame}

\begin{frame}[wide]{The Model’s Prediction for Per Capita GDP}
	\begin{itemize}
		\item If the model were successful in explaining incomes, countries should lie close to the solid 45-degree line
		\item Instead, the model predicts that most countries should be substantially richer than they are		
	\end{itemize}
	\begin{figure}
		\centering
		\includegraphics[width=0.5\linewidth]{"Figures/The Model’s Prediction for Per Capita GDP_figure"}
	\end{figure}
	
\end{frame}

%\begin{frame}[wide]{Case Study: Why Doesn't Capital Flow from Rich to Poor Countries?}
%	If MPK is higher in poor countries with low K, why doesn't capital flow to those countries?
%	\begin{itemize}
%		\item Caselli and Feyrer use data on GDP, K, and the shape of the production function to measure the MPK directly for many countries around the world
%		\item They find that MPK is quite similar across a range of countries
%		\begin{itemize}
%			\item MPK in rich countries is slightly greater than MPK in poor countries.
%			\item 8.4 percent in rich countries versus 6.9 percent in poor
%		\end{itemize}
%		\item The puzzle, then, is not why capital fails to flow to poor countries
%		\item The puzzle is, rather, why the MPK in poor countries is not much higher, given that they have so little capital (e.g., economic policy, education, misallocation, etc)
%	
%		
%	\end{itemize}
%\end{frame}

\begin{frame}[wide]{Productivity Differences: Improving the Fit of the Model}
	\begin{itemize}
		\item 	The productivity parameter measures how efficiently countries are using their factor inputs -- Total factor productivity (TFP)
		\begin{itemize}
			\item A lower level of TFP implies that workers produce less output for any given level of capital per person
		\end{itemize}
		\item If TFP is not equal to one for every countries, it would lead to a better model
		\item But, we have data limitation in that there is no real direct measure of TFP
		\item We have to exploit the fact that we have data on every term in the equation other than TFP
		\begin{itemize}
			\item That's why TFP is sometime refer as Solow's residual
		\end{itemize}
		\item Therefore, we can calculate the level of TFP for each country that would be needed to make the equation hold exactly 
	\end{itemize}
\end{frame}

\begin{frame}[wide]{Measuring TFP so the Model Fits Exactly}
	\begin{itemize}
		\item 	In order for our model to match the data,	poor countries must be very inefficient in production, or 	they must have low TFP
		\item For example, 	$\text{Italy} = \frac{y}{k^{1/3}} = 0.68/0.93 = 0.731 $
		
	\end{itemize}
	\begin{figure}
		\centering
		\includegraphics[width=0.6\linewidth]{"Figures/Measuring TFP so the Model Fits Exactly"}
	\end{figure}
	
\end{frame}

\begin{frame}[wide]{Measuring TFP so the Model Fits Exactly}
	\begin{itemize}
		\item Notice that TFP and per capita GDP are highly correlated
		\item Also, TFP varies substantially across countries
		\item This is the motivation for examining economic growth and understanding differences between rich and poor countries
	\end{itemize}
	\begin{figure}
		\centering
		\includegraphics[width=0.5\linewidth]{"Figures/Measuring TFP so the Model Fits Exactly—2"}
	\end{figure}
	
\end{frame}

\begin{frame}[wide]{Understanding TFP Differences}
	Based on the simple production function:
	\vspace{2mm}
	\begin{itemize}
		\item 	Output differences between the richest and poorest countries?
		\begin{itemize}
			\item Differences in capital per person explain about one-third of the difference
			\item TFP explains the remaining two-thirds
		\end{itemize}
		\item Thus, rich countries are rich because
		\begin{itemize}
			\item they have more capital per person
			\item more importantly, they use labour and capital more efficiently
		\end{itemize}
		\item Why are some countries more efficient at using capital and labour?
	\end{itemize}
\end{frame}

\begin{frame}[wide]{Understanding TFP Differences - Human Capital}
	\begin{itemize}
		\item Human Capital: the stock of knowledge and skills that make people more productive members of society
		\begin{itemize}
			\item e.g., Education, training, health, etc
		\end{itemize}		
		\item A lot of time researchers measure it through \textit{returns to education}: the value of the increase in wages from additional schooling
		\item There are significant returns to education within a country
		\begin{itemize}
			\item In the United States each year of education seems to increase future wages by 7 percent
			\begin{itemize}
				\item A four-year education may raise wages by about 28 percent (over entire lifetime)
			\end{itemize}
			\item In developing countries, returns can be even higher—up to 10 percent or even 13 percent per year
			\item The typical student learning the basic skills associated with literacy and arithmetic may have higher returns than a college education.
		\end{itemize}
		\item Accounting for human capital can reduces the residual 
	\end{itemize}
\end{frame}

\begin{frame}[wide]{Understanding TFP Differences - Technology}
	\begin{itemize}
		\item 	Richer countries may use more modern and efficient technologies than poor countries
		\begin{itemize}
			\item Increases productivity parameter
		\end{itemize}
		\item But, how does the technology diffuse across countries?
	\end{itemize}
\end{frame}

\begin{frame}[wide]{Understanding TFP Differences - Institutions}
	\begin{itemize}
		\item Well-defined institutions and laws create a climate for economic growth
		\item Institutions are in place to foster human capital and technological growth
		\begin{itemize}
			\item Property rights
			\item The rule of law
			\item Government systems
			\item Contract enforcement
		\end{itemize}
		\item At the time of the collapse of the Berlin Wall in 1989, the standards of living in East and West Germany were substantially different
		\item Difference between Civil law and Common law can lead to difference in economic development
	\end{itemize}
\end{frame}

\begin{frame}[wide]{Understanding TFP Differences - Misallocation}
	\begin{itemize}
		\item Misallocation
		\begin{itemize}
			\item Resources not being put to their best use
		\end{itemize}
		\item Examples
		\begin{itemize}
			\item Inefficiency of state-run resources
			\item Political interference
		\end{itemize}
		\item Inefficient land policy: small land size -> farm cannot be mechanize 
%		\begin{itemize}
%			\item Restuccia has made significant contribution to the resource allocations literature
%		\end{itemize}
	\end{itemize}
\end{frame}

\begin{frame}[wide]{Evaluating the Production Model}
	\begin{itemize}
		\item Per capita GDP is higher if capital per person is higher and if factors are used more efficiently.
		\item Constant returns to scale imply that output per person can be written as a function of capital per person
		\item Capital per person is subject to strong diminishing returns because the exponent is much less than one
	\end{itemize}
\end{frame}

\begin{frame}[wide]{Weaknesses of the Model}
	\begin{itemize}
		\item The production model has partial empirical success
		\item In the absence of TFP, the production model predicts only a small differences in income
		\item The model provided little explanation for why countries exhibit different TFP levels and different capital
		\item Why do economies grow over time?
		\item This is what we explore in the next two chapters
	\end{itemize}
\end{frame}


%\begin{frame}[wide]{Measuring the State of the Economy}
%	\begin{columns} 
%		% Column 1
%		\begin{column}{.5\textwidth}
%			\begin{table}[htbp]
%				\centering
%				\begin{tabular}{lrrrr}
%					& 2005  & 2006  & 2007  & 2008 \\
%					\hline
%					GDP   & 12.6  & 13.4  & 14.0    & 14.3 \\
%					(in trillions of \$) &       &       &       &  \\
%					Growth rate GDP &       & 0.063 & 0.045 & 0.021 \\
%					\hline
%				\end{tabular}%
%				\label{tab:addlabel}%
%			\end{table}%
%		\end{column}
%		% Column 2    
%		\begin{column}{.5\textwidth}			
%
%			
%		\end{column}%
%	\end{columns}
%\end{frame}

\begin{frame}[wide]{Next Lecture}
	\begin{itemize}
		\item Next lecture will start Ch.5
%		\item Group assignment is available on Crowdmark 
%		\begin{itemize}
%			\item Group up to four students
%			\item Once you started, you will have only 3 hours to complete the assignment
%			\item Leave enough time for pdf/image upload
%		\end{itemize}
	\end{itemize}
\end{frame}

\end{document}